\documentclass[10pt,a4paper]{amsart}
\usepackage[margin=19mm]{geometry}
\usepackage{amsmath,amssymb,mathtools}
\usepackage[scr=rsfs,cal=euler]{mathalfa}
\usepackage{xfrac}
\usepackage[unicode,hidelinks,bookmarks=false]{hyperref}
\newcommand{\ie}{i.e.\ }
\newcommand{\cf}{cf.\ }
\newcommand{\resp}{respectively\ }
\newcommand{\Subsection}{\subsection}
\providecommand{\nobreakdash}{\nobreak\hbox{-}}
\setlength{\emergencystretch}{2em}
\multlinegap0pt
\usepackage{bm}
\usepackage{bbm}
\usepackage{dsfont}
\usepackage{xspace}
\usepackage{cancel}
\usepackage{mathtools}
\usepackage{amsthm}
\makeatletter
\@ifundefined{lemma}{\newtheorem{lemma}{Lemma}}{}
\makeatother
\title[Thin Trees via $k$-Respecting Cut Identities]{Thin Trees via $k$-Respecting Cut Identities}
\author{Mohit Daga}
\date{Source: arXiv:2510.12050v1 (CC BY 4.0)}
\begin{document}
\maketitle

\noindent\emph{Source and licence.} Thin Trees via $k$-Respecting Cut Identities. Version arXiv:2510.12050v1 (CC BY 4.0), distributed under the Creative Commons Attribution 4.0 International licence, as recorded in the arXiv licence metadata for this exact version and shown on the abstract page https://arxiv.org/abs/2510.12050v1. Full rights notice: licenses/arxiv-2510.12050-CC-BY-4.0.txt.\par\smallskip

\noindent\textit{Editorial scope.} Counted unit: the Lemma whose source label is \texttt{lem:krespect-eval} in the article (source0.tex lines 782--803, source label \texttt{lem:krespect-eval}), delivered together with the article's own complete original proof of it (source0.tex lines 805--828, one \texttt{proof} environment of 24 lines). No mathematics is written, repaired or re-derived: statement and proof are the article's own text line for line. Required context retained uncounted: none. Every definition, display and label of those lines is retained unchanged. Omitted from this package and not redistributed: 1--781, 804--804, 829--1858. Nothing inside a retained excerpt is omitted or altered: every retained body is delivered exactly as the article's lines are written, so no omission receipt of any kind accompanies this package. Citations: the retained excerpts carry no citation command at all, so no bibliography entry of the article is transcribed and no reference is redistributed. Reference adaptation: none was needed and none was made; every reference inside the delivered unit resolves inside it, so no unresolved reference or citation is produced. Printed slips kept literal: no printed word, symbol, hypothesis or number of the counted statement or of its proof was altered. Layout and class: the article's own class is \texttt{article} with packages including tikz, pgfmath, hyperref, geometry, inputenc, babel, ulem, amsmath, amssymb, cases, amsthm, thmtools, enumitem, mdframed; the wrapper is the lane's pinned \texttt{amsart} editorial wrapper at 10pt on A4, which restates the theorem-like environments the retained text needs and adds the article's own macro definitions that the retained text needs (none), with extra wrapper packages (bm, bbm, dsfont, xspace, cancel, mathtools). The delivered numbering is the unit's own. Assets: no retained span contains a figure environment, a graphics-inclusion command, a PDF-page inclusion, a table or a TikZ picture; no image file is copied into this package, the package declares no asset dependency and no image file is redistributed with it. Licence: version arXiv:2510.12050v1 of this article is distributed under the Creative Commons Attribution 4.0 International licence, as recorded in the arXiv licence metadata for this exact version and shown on the abstract page https://arxiv.org/abs/2510.12050v1. A full rights notice is delivered at licenses/arxiv-2510.12050-CC-BY-4.0.txt. Characters: every retained excerpt is pure ASCII, so the delivered engine, XeLaTeX with its default Latin Modern text font, reports no missing character.
\begin{lemma}[Exact $k$--respecting cut evaluation from pairwise data]
\label{lem:krespect-eval}
Fix a rooted spanning tree $T$ and let $A\subseteq V$ be a cut
such that $|E(T)\cap \delta(A)|=k$.
Let $S=\{v_1,\dots,v_k\}$ be the child endpoints of the $k$ tree edges
crossed by $\delta(A)$, so that
\[
A \;=\; \bigoplus_{i=1}^k D_T(v_i).
\]
Then
\begin{equation}
\label{eq:inclusion-exclusion}
w(\delta(A)) \;=\;
\sum_{\ell=1}^k (-1)^{\ell-1} 2^{\ell-1}
  \sum_{\substack{I\subseteq [k]\\ |I|=\ell}}
  w\!\left( \bigcap_{i\in I} \delta(D_T(v_i)) \right).
\end{equation}
Moreover each term on the right--hand side is determined solely by the pairwise
values $\sigma(u,v)$ and the ancestor--descendant relations in $T$.
Hence, after $O(n^2)$ preprocessing of all $\sigma(u,v)$ values,
the quantity $w(\delta(A))$ can be evaluated in $O(k^2)$ time.
\end{lemma}

\begin{proof}
The expression of $A$ as a symmetric difference of $k$ subtrees
follows from rooting $T$: membership in $A$ flips whenever a
crossing edge of $\delta(A)$ is traversed, so the shore $A$
is obtained exactly by toggling the descendant sets of the $k$ child
vertices $v_1,\dots,v_k$.

Consider any edge $e=\{x,y\}\in E$.
The indicator $[e\in \delta(A)]$ equals the parity of the number of
subtrees among $\{D_T(v_i)\}$ that separate $x$ and $y$.
In other words,
\[
[e\in \delta(A)] \;=\; \bigoplus_{i=1}^k [e\in \delta(D_T(v_i))].
\]
Expanding the parity by inclusion--exclusion gives exactly the formula
\eqref{eq:inclusion-exclusion}, with coefficient $2^{\ell-1}$ for the
$\ell$-wise intersections: an edge is cut by an odd number of subtrees
iff it belongs to an odd number of the sets $\delta(D_T(v_i))$,
which is captured by alternating signs and powers of two.

It remains to argue that each term, in $\sum_{\substack{I\subseteq [k]\\ |I|=\ell\geq2}}
  w\!\left( \bigcap_{i\in I} \delta(D_T(v_i)) \right)$ of Eq \ref{eq:inclusion-exclusion}, for example,
$w(\delta(D_T(v_i))\cap \delta(D_T(v_j)))$ and other higher order terms are determined by the pairwise values $\sigma(\cdot,\cdot)$.
\end{proof}

\end{document}
